%% AASTeX v6.31 manuscript skeleton. Target: The Planetary Science Journal (alt: AJ).
%% Compile on Overleaf (AAS template) or with a local TeX Live:  latexmk -pdf ms.tex
%% Every number in this file is a macro defined in numbers.tex, which
%% scripts/make_numbers_tex.py regenerates from paper/results/*.json.
%% Lines starting with "%% TODO" are writing tasks, not content.
\documentclass[twocolumn,linenumbers]{aastex631}

\usepackage{amsmath}
% generated by scripts/make_numbers_tex.py -- do not edit
\newcommand{\NbbSixteen}{31}
\newcommand{\KbbOrbits}{2.5}
\newcommand{\KbbNondet}{1.7}
\newcommand{\KselMarg}{5.6}
\newcommand{\KselMargCont}{4.2}
\newcommand{\lnKpeak}{5.9}
\newcommand{\YearPeak}{2013}
\newcommand{\KpreTwoThousandSixteen}{17}
\newcommand{\NtrainObjects}{23}
\newcommand{\NnewObjects}{8}
\newcommand{\lnKnew}{-1.87}
\newcommand{\KnewObjects}{0.15}
\newcommand{\PnewUnderPlanet}{3.1\%}
\newcommand{\PnewUnderNull}{25\%}
\newcommand{\NKhundredLSST}{47--95}
\newcommand{\NKhundredCurrent}{63--133}
\newcommand{\YieldLSST}{233}
\newcommand{\YieldLSSTrange}{134--409}
\newcommand{\YieldYears}{3--8 years}
\newcommand{\YieldYearsMedian}{2--5 years}
\newcommand{\PredictiveRows}{2013 & baseline & 11 & 20 & -4.89 & 0.000 & 0.35 \\
2014 & baseline & 16 & 15 & -4.24 & 0.001 & 0.19 \\
2015 & baseline & 23 & 8 & -1.87 & 0.031 & 0.25 \\
2015 & $a>50$\,au footprint & 23 & 8 & -1.64 & 0.056 & 0.41 \\
2016 & baseline & 25 & 6 & -2.37 & 0.009 & 0.09 \\}
\newcommand{\NbodySummary}{\textbf{??}}


\newcommand{\Mearth}{M_\oplus}
\newcommand{\vpi}{\varpi}
\newcommand{\Kbf}{\mathcal{K}}

\shorttitle{A prequential test of Planet Nine}
\shortauthors{B. Thelen}

\begin{document}

\title{Planet Nine's orbital evidence, tested out of sample: the Bayes factor has fallen since 2016, and when Rubin will decide}

%% TODO author list, affiliations, ORCID. The submitting author stands behind every number.
\author[0000-0000-0000-0000]{Benjamin Thelen}
\affiliation{Independent researcher}

\begin{abstract}
%% 250 words max. Structure: question -> method -> three results -> forecast.
Whether the clustered orbits of the most distant trans-Neptunian objects (ETNOs) require an undiscovered planet has been argued for a decade with significance tests whose answers depend on how survey bias is modelled.
We compute the Bayes factor $\Kbf$ between a planet that confines part of the population and an isotropic population, weighting every object by its own selection function, and treat the selection model itself as an uncertain nuisance.
For the \NbbSixteen{} known ETNOs with $a>250$\,au and $q>30$\,au we find $\Kbf=\KbbNondet$ with the planet's non-detection included, \KbbOrbits{} from the orbits alone, and at most \KselMarg{} after marginalising over selection models: ``barely worth mentioning'' to ``substantial''.
Ordered by discovery date, the evidence peaked at $\ln\Kbf=\lnKpeak$ with the objects known by \YearPeak{} and has fallen almost monotonically since.
This makes a prequential test possible: the planet inferred from the \NtrainObjects{} objects found by 2015 predicts the \NnewObjects{} later ones worse than isotropy ($\Kbf_\mathrm{new}=\KnewObjects$), an outcome that occurs in only \PnewUnderPlanet{} of posterior-predictive skies but is typical of skies without a planet. Every split year from 2013 to 2016 gives $p\le0.03$.
Gigayear $N$-body integrations of three published planet configurations confirm the anti-aligned, partially confined population that the likelihood assumes (\NbodySummary).
Finally we model the Rubin Observatory LSST selection function and forecast the decision: about \NKhundredLSST{} LSST-discovered ETNOs are needed for $\Kbf=100$ if the planet the current data imply is real, which an OSSOS-anchored population model puts at \YieldYears{} of survey operations.
We publish the decision thresholds, the full pipeline and the predicted $\Kbf$ distributions in advance, as a pre-registered test for the LSST era.
\end{abstract}

\keywords{Trans-Neptunian objects (1705) --- Kuiper belt (893) --- Planetary dynamics (2173) --- Bayesian statistics (1900) --- Sky surveys (1464)}

%% ------------------------------------------------------------------
\section{Introduction}\label{sec:intro}
%% TODO ~1 page. Paragraph plan:
%% (1) The clustering claim: Trujillo & Sheppard 2014; Batygin & Brown 2016; Brown & Batygin 2016, 2019, 2021; review Batygin et al. 2019.
%% (2) The bias counter-argument: Shankman et al. 2017 (OSSOS); Brown 2017 (bias model); Napier et al. 2021 (DES+OSSOS+ST, no evidence); Bernardinelli et al. 2020/2022.
%% (3) The newest assessment: Siraj, Chyba & Tremaine 2026 (conditional likelihood, 2.7 -> 1.9 sigma as the sample grew 21 -> 25). State precisely how this work differs:
%%     (a) Bayes factor against an explicit planet model rather than a significance against isotropy;
%%     (b) the decline treated as a calibrated out-of-sample test with posterior-predictive simulations;
%%     (c) selection-model uncertainty marginalised, not fixed;
%%     (d) a quantitative, pre-registered LSST decision forecast. No published ETNO-specific LSST yield exists (Kurlander et al. 2025 give ~37,000 TNOs in total).
%% (4) Alternatives that also produce clustering or warps: self-gravitating disk (Madigan & McCourt 2016; Sefilian & Touma 2019; Zderic & Madigan 2020), stellar flyby (Pfalzner 2026), "Planet Y" mean-plane warp (Siraj, Chyba & Tremaine 2025), Kuiper-belt planet (Lykawka & Ito 2023). Our H1 is the Planet Nine family only.
%% (5) Roadmap of the paper.
The orbits of the most distant trans-Neptunian objects are clustered in longitude of perihelion and in orbital pole \citep{Trujillo2014,Batygin2016}. \dots

%% ------------------------------------------------------------------
\section{Data}\label{sec:data}
%% Source: MPC Distant.txt export with orbits to 2026-01-22 (6,536 orbits), arc >= 1 yr, numeric U.
%% Table 1: the three samples (bb16 N=31; ts14 N=62; detached N=19).
%% Table 2 (appendix A): the 31 bb16 objects with a, q, i, H, varpi, Omega, discovery year, discovering survey, per-object ln K contribution.
%% TODO before submission: refresh to the latest MPC/JPL orbits (this environment could not reach them) and record the snapshot DOI on Zenodo.
%% Discovery year = provisional-designation year. Caveat: announcement can lag discovery by years (e.g. 2017 OF201, announced 2025; Cheng et al. 2025). The prequential split is repeated with the BB16 training set (6 objects) as a check.

\begin{deluxetable}{lccc}
\tablecaption{ETNO samples (MPC snapshot 2026-01-22)\label{tab:samples}}
\tablehead{\colhead{Sample} & \colhead{Definition} & \colhead{$N$} & \colhead{Use}}
\startdata
bb16 & $a>250$\,au, $q>30$\,au & \NbbSixteen & baseline \\
detached & $a>250$\,au, $q>40$\,au & 19 & robustness \\
ts14 & $a>150$\,au, $q>30$\,au & 62 & regime check \\
\enddata
\end{deluxetable}

%% ------------------------------------------------------------------
\section{Method}\label{sec:method}

\subsection{Per-object selection functions}\label{sec:selection}
%% Empirical footprint (Brown 2017 style): 5,396 distant discoveries mark surveyed sky to a depth; S_j(varpi, Omega) on a 72x72 grid.
%% Characterised surveys: OSSOS v11 family (Bannister et al. 2018) pointings and efficiencies for the 9 ETNOs they found.
%% Independent implementation (planetix.bias) as a cross-check.

\subsection{Likelihood and planet model}\label{sec:likelihood}
%% Eq. 1: p(x_j | H) = f(x_j) S_j(x_j) / \int f S_j dx.
%% Eq. 2: mixture f = (1-F_j) + F_j vM(varpi; varpi9+180, kappa_varpi) Fisher(pole; pole9, kappa_pole).
%% Eq. 3: F_j sigmoid in ln a around a_crit(m9, a9, e9); a_crit normalisation and secular check (planetix.dynamics).
%% Table 3: priors.

\subsection{Selection model as a nuisance parameter}\label{sec:selmarg}
%% New. Absolute H0 likelihood ln p(D | H0, S_k) = sum_j ln[S_jk(x_j) n_cells]; posterior weight of each selection model
%% w_k \propto p(D | H0, S_k) (equal model priors), and K_marg = sum_k w_k Z0_k K_k / sum_k w_k Z0_k.
%% Continuous family: geometric interpolation S(alpha) \propto S_r2^(1-alpha) S_alt^alpha, alpha ~ U(0,1).
\begin{equation}
\Kbf_\mathrm{marg} = \frac{\int p(\alpha)\,Z_0(\alpha)\,\Kbf(\alpha)\,d\alpha}{\int p(\alpha)\,Z_0(\alpha)\,d\alpha}.
\end{equation}

\subsection{Non-detection of the planet}\label{sec:nondet}
%% Pan-STARRS1/ZTF coverage proxy; Chen & Kipping 2017 mass-radius. Excludes 64% vs 78% published (Brown et al. 2024): conservative toward H1.

\subsection{Evidence, calibration and the prequential decomposition}\label{sec:prequential}
%% Nested sampling (dynesty; Speagle 2020). Calibration catalogues under H0, H0 with wrong selection, and H1.
%% Prequential identity (Dawid 1984): ln K(D_1..N) = sum_n ln[p(x_n | x_<n, H1) / p(x_n | H0)], so ln K(all) - ln K(train) is the evidence carried by the new objects given what the old ones taught.
%% Test statistic: K_new = p(x_new | x_old, H1) / p(x_new | H0), estimated with training-posterior samples;
%% reference distributions from posterior-predictive (planet real, as the old objects imply) and from H0 skies.

\subsection{$N$-body forward model}\label{sec:nbody}
%% New. REBOUND WHFast; Sun + planet + 1000 test particles; giants as orbit-averaged J2 applied as exact secular rotations;
%% Neptune as a q < 30 au sink; 1 Gyr; three planets (BB16, reference, BB21).
%% Output: confined fraction, anti-alignment and pole offset vs a; compared with the mixture model's F(a), kappa_varpi, kappa_pole.

%% ------------------------------------------------------------------
\section{Results}\label{sec:results}

\subsection{The Bayes factor today}\label{sec:bf}
%% Table 4: all variants (main, nondet, detached, ts14, bias footprint, OSSOS-characterised, selection-marginalised).
%% Figure 1 (fig1_orientations): orientations over the mean selection function.
%% Figure 2 (fig2_bayes_factors): ln K for every variant.
For the baseline sample $\Kbf=\KbbOrbits$ from the orbits alone and \KbbNondet{} with the non-detection term.

\subsection{Evidence history}\label{sec:history}
%% Figure 6 (fig6_history): ln K vs discovery year for bb16 and ts14, with sample size on a twin axis and the Jeffreys bands.
%% Text: peak value, year of peak, decline since.
The evidence peaked at $\ln\Kbf=\lnKpeak$ ($\Kbf\approx360$, ``decisive'') with the 11 objects known by \YearPeak{}, was $\KpreTwoThousandSixteen$ with the objects known by 2015, and has declined at almost every step since (Figure~6). The $a>150$\,au sample shows the same shape with a later peak (2016).
%% TODO one paragraph: which surveys found the early objects (Sheppard-Trujillo, OSSOS) vs the later ones; regression to the mean after a hypothesis chosen on the data.

\subsection{An out-of-sample test of the pre-2016 planet}\label{sec:oos}
%% Figure 7 (fig7_predictive): distribution of ln K_new under the trained planet (posterior predictive) and under H0, observed value marked.
%% Table 6 (appendix): per-object ln K_new contributions of the post-2015 objects (predictive_2015.json).
\begin{deluxetable}{lcccccc}
\tablecaption{Out-of-sample tests: the planet inferred from objects discovered up to the split year, scored on the later objects\label{tab:oos}}
\tablehead{\colhead{Split} & \colhead{Selection} & \colhead{$N_\mathrm{train}$} & \colhead{$N_\mathrm{test}$} & \colhead{$\ln\Kbf_\mathrm{new}$} & \colhead{$p$ (trained planet)} & \colhead{$P$ (no planet, $\le$ obs)}}
\startdata
\PredictiveRows
\enddata
\tablecomments{$p$: fraction of 1,000 posterior-predictive skies with $\ln\Kbf_\mathrm{new}$ at or below the observed value. Last column: the same fraction for skies without a planet.}
\end{deluxetable}
Trained on the \NtrainObjects{} objects discovered by 2015, the planet hypothesis predicted the orientations of the \NnewObjects{} later objects worse than isotropy: $\ln\Kbf_\mathrm{new}=\lnKnew$. Under the planet the training set implies, $\ln\Kbf_\mathrm{new}$ is this low in \PnewUnderPlanet{} of simulated skies; under no planet it is this low or lower in \PnewUnderNull{}.

\subsection{Selection-marginalised evidence}\label{sec:selres}
%% Figure 8: ln K(alpha) and Z0(alpha); table of discrete models with weights.
Averaging over selection models with weights set by the data gives $\Kbf=\KselMarg$ (discrete set) and \KselMargCont{} (continuous family). The data favour the footprint built only from $a>50$\,au discoveries (posterior weight 0.99 in the discrete set), which raises $\Kbf$; the conclusion that the orbits give at most substantial evidence does not depend on the selection model.

\subsection{$N$-body check of the confinement model}\label{sec:nbodyres}
%% Figure 9 (fig9_nbody): mean resultant length of varpi - varpi9 and pole offset vs a for each planet, with the model's F(a) overplotted.
%% Key numbers: confined fraction at a >> a_crit, anti-alignment centre, kappa_varpi; whether the analytic a_crit matches.
\NbodySummary.

\subsection{Planet parameters if the planet is adopted}\label{sec:posterior}
%% Figure 4 (fig4_posterior). Orientation constrained (varpi9, Omega9), mass and distance prior-dominated along m9 \propto a9^3.

%% ------------------------------------------------------------------
\section{Forecast: when Rubin decides}\label{sec:forecast}

\subsection{An LSST selection model}\label{sec:lsstsel}
%% WFD dec -70..+12.5, |b|>15; NES |beta|<10 to dec +40; V_lim 24.5 (24.0-25.5 range); Ivezic et al. 2019.
%% Figure 10: S_j for three objects, current footprint vs LSST.

\subsection{Information rate and decision sample size}\label{sec:info}
%% E[ln K(N)] ~ N <D> - c ln N - c0, with <D> the per-object KL divergence (exact on the grid), c, c0 fitted to injections.
%% Validated with end-to-end simulated LSST catalogues (lsstsim).
%% Figure 11 (fig11_decision): N needed for K=100 over (f_max, kappa_varpi), current vs LSST selection; observed posterior overplotted.
If the planet implied by today's data is real, $\Kbf=100$ needs about \NKhundredLSST{} ETNOs found by LSST (\NKhundredCurrent{} with today's mixture of surveys).

\subsection{Yield and timeline}\label{sec:yield}
%% OSSOS-anchored: N_LSST = N_OSSOS,obs <P_LSST>/<P_OSSOS>, Poisson (Jeffreys) interval on N_OSSOS,obs = 3, population-model variants.
LSST should find \YieldLSST{} ETNOs with $a>250$\,au over ten years (68\% range \YieldLSSTrange{}; 114--278 across population models; 100--525 for $V_\mathrm{lim}=24.0$--$25.0$). LSST detects them near perihelion: restricting detection to $r<100$\,au, where the moving-object pipeline is efficient, removes none. If discoveries accrue evenly after a one-year orbit-arc delay, the median sky reaches $\Kbf=100$ after \YieldYearsMedian{} and 90\% of skies after \YieldYears{}.

\subsection{Pre-registered decision rule}\label{sec:prereg}
%% Fixed now, before LSST ETNO orbits exist: sample = LSST-discovered objects with a>250, q>30, arc >= 2 oppositions;
%% pipeline = this paper's code at tag v1.0 (Zenodo DOI); selection = LSST survey simulator (Sorcha; Merritt et al. 2025) when public, else Sec. 5.1 model;
%% decision: K >= 100 -> evidence for the planet; K <= 1/30 -> evidence against; otherwise inconclusive, re-evaluate at each data release.
%% Predicted K distribution vs N under H0 and implied H1 tabulated in Table 6.

%% ------------------------------------------------------------------
\section{Discussion}\label{sec:discussion}
%% 6.1 Comparison with Napier et al. 2021, Brown & Batygin 2021, Siraj et al. 2026 (agreement of direction and size).
%% 6.2 Why the evidence fell: regression to the mean after a hypothesis chosen on the data (a winner's-curse effect); quantify with the calibration.
%% 6.3 The ts14 signal (a > 150 au): needs a closer heavier planet; Neptune-driven? look-elsewhere p = 0.064.
%% 6.4 Limitations: empirical selection proxy; announcement dates; orientation-only likelihood; no alternative H1; nondetection model coarse.

\section{Conclusions}\label{sec:conclusions}
%% Four numbered findings + the pre-registered test.

\begin{acknowledgments}
This work used the Minor Planet Center's database and the OSSOS survey data release.
\end{acknowledgments}

\software{numpy, scipy, pandas, dynesty \citep{Speagle2020}, REBOUND \citep{Rein2012,Rein2015}, matplotlib}

\appendix
\section{Reproducibility}\label{app:repro}
%% Code: github.com/1bnjmn3/Planet-IX (oct26-edits/), Zenodo DOI TBD. Commands: scripts/infer_planet9.py, scripts/forecast_and_history.py, scripts/nbody_population.py, scripts/make_paper_figures.py, scripts/make_numbers_tex.py.
\section{Per-object data}\label{app:objects}

\bibliography{refs}{}
\bibliographystyle{aasjournal}

\end{document}
